\chapter{Where Humanvoice fits}
\label{ch:field}

Take one source file and follow it to a reader's decision. Which existing
components can keep that path inspectable? ``Humanize prose'' is not a validated
package objective; searching for it would put the surface before the reader's
task. The narrower question gives us a basis for comparison.

Figure~\ref{fig:capability-map} gives the organizing answer. Mature components
surround a record that no current package owns; that record is the proposed
product boundary.

\hvFigureCapabilityMap

\section{The available components}

\subsection{Representing a technical document}

When an author changes a sentence beside an equation, the system must know
which source span produced each object. Pandoc exposes a typed document
abstract syntax tree to filters \citep{pandoc2026filters}. pylatexenc parses
LaTeX into nodes and allows project-specific macro and environment handling
\citep{pylatexenc2026}. plasTeX expands macros into a document model, while
tree-sitter-latex offers incremental syntax trees; LaTeXML is a stronger
conversion path when macro expansion and semantic output are required
\citep{plastex2026,treesitterlatex2026,latexml2026}. The smaller
\texttt{latexlint} utility is a useful reference for format-aware checks, but
its role remains a benchmark candidate until its location and protected-object
behavior are tested on this repository's macros \citep{latexlint2026}.

Two additional candidates close a language and tolerance gap in the original
shortlist. unified-latex exposes a JavaScript/TypeScript abstract syntax tree
and transformation utilities in the unified ecosystem
\citep{unifiedlatex2026}. TexSoup offers a compact, tolerant Python interface
that is attractive for prototypes and deliberately messy fixtures
\citep{texsoup2026}. Neither documentation set establishes source fidelity on
the custom macros in this repository; both now belong in the same locked
parser benchmark.

No single parser should be declared correct from a README. The MVP will freeze
fixtures containing equations, labels, citations, comments, custom macros,
tables, malformed input, and round-trip locations. Lightweight and full-model
parser candidates will compete on those fixtures. The winner is the one that preserves the objects the
protected comparison needs, not the one with the most convenient marketing
description.

\subsection{Finding observable problems}

Vale supplies configurable, markup-aware prose rules
\citep{vale2026docs}. LanguageTool and LTeX+ provide grammar and language
diagnostics that can run locally or through an editor protocol
\citep{languagetool2026dev,ltexplus2026}. textlint offers a pluggable rule
interface and machine-readable output \citep{textlint2026docs}. These are useful adapters for local
observations: a repeated opener, a banned internal identifier, a grammar
violation, or a sentence that needs an author question.

They do not supply a reader-acceptance layer. A tool can report that a phrase
matches a rule; it cannot establish that the phrase is appropriate for this
audience, or that removing it improves the argument. Humanvoice must keep the
rule output separate from the author decision.

\subsection{Comparing and checking sources}

latexdiff makes a visual LaTeX diff \citep{latexdiff2026}, but cannot prove
that equations or citations survived. GROBID extracts PDF references
\citep{grobid2026}; LaTeX should use its source bibliography plus metadata. It
separates citation identity from sentence support; substantive changes need
human verification.

\subsection{Models and evaluation infrastructure}

llama.cpp makes local inference practical on supported hardware
\citep{llamacpp2026}. Inspect AI provides structured datasets, solvers, scorers,
and logs for model evaluations \citep{inspectai2026}. Both are candidates for
privacy and reproducibility, not evidence that a local model gives good
editorial advice. The MVP can run without either: deterministic planning,
bounded drafting, pre-human diagnostics, and a fail-closed release gate are the
core path, while model assistance remains a replaceable adapter with an explicit
disclosure record. A model may write or criticize a unit, but it cannot release
its own work.

\section{The benchmark's decisions}

The shortlist matters only when it changes a build decision. In the matrix,
``candidate'' means that a component has a plausible job; it does not mean that
the component works on Humanvoice documents. The current audit has smoke-level
results for only some local tools, so held-out fixtures and human labels still
stand between every candidate and adoption.

\begin{table}[H]
\centering
\small
\begin{tabularx}{\textwidth}{@{}p{0.19\textwidth}p{0.19\textwidth}L p{0.22\textwidth}@{}}
\toprule
Need & Candidate components & Decision criterion & First-release role \\
\midrule
Source map & Pandoc, pylatexenc, unified-latex, TexSoup, plasTeX,
tree-sitter-latex & Protected-region
and line-location fixtures; malformed-source recovery; round-trip fidelity. &
Select one primary adapter and retain one fallback. \\
Prose diagnostics & Vale, LTeX+, LanguageTool, textlint & Precision and burden
on held-out technical documents; markup exclusions; stable machine output. &
Use replaceable deterministic adapters. \\
Visual comparison & latexdiff & Review readability and source-link fidelity. &
Supplement, never replace, protected invariants. \\
Citation identity & Source bibliography plus Crossref or publisher metadata;
GROBID for PDF ingestion & Correct classification of valid, wrong, duplicate,
and unresolved identities. & Block unresolved load-bearing identities; queue
support review. \\
Model assistance & Optional local runtime and Inspect AI & Calibration against
human labels, privacy, order/length sensitivity, and abstention. & Optional
screening only; no release authority. \\
\bottomrule
\end{tabularx}
\caption{Software choices are hypotheses tested against the product contract,
not a popularity ranking.}
\label{tab:software-shortlist}
\end{table}

The benchmark will report runtime, failure modes, license and maintenance
status, privacy boundary, and human review burden alongside precision. A
component that is accurate but produces an unmanageable queue is not suitable
for the workflow. A component that is easy to install but cannot preserve
protected spans is not suitable for technical writing.

\section{The existing project is an advantage only if it stays subordinate}

The repository already contains three assets that shorten the path to a pilot:
recorded cases of reader rejection and repair, a writing policy that places
meaning before regularity, and scripts that can render and inspect LaTeX
documents. Those assets supply fixtures and vocabulary. They do not establish
that a diagnostic improves a reader's decision. The new code should turn them
into a small corpus and a reproducible review packet, not into another layer of
rules.

\section{The narrow advantage}

The nearest alternatives are not package names. They are ways a team might
spend the next month:

\begin{description}
\item[Generic LLM editing.] Fast and flexible, but it leaves source protection,
  provenance, and reader acceptance to local discipline.
\item[Enterprise writing assistance.] Convenient and integrated, but often
  optimized for general business prose and remote processing rather than
  equations, citations, and confidential research.
\item[Rules plus manual review.] Transparent and controllable, but it leaves
  the author to connect a finding to a revision and a reader outcome.
\item[Humanvoice.] A deliberately narrow connection between those activities,
  with the human reader retained as the acceptance authority.
\end{description}

The wedge is not a superior language model. It is a durable record of why a
change was proposed, what meaning-bearing objects survived, and whether the
intended reader could act. That record can survive a change of parser, linter,
or model provider. It is also the reason the product can be evaluated against
the incumbent bundle rather than against an imaginary perfect editor.

\begin{boundarybox}{Adoption rule}
No package earns a permanent place because it is popular, current, or easy to
install. It earns a place when it passes the protected-source fixtures, has a
clear privacy and maintenance story, and changes a measured workflow outcome
on held-out documents. Until then it remains a replaceable candidate.
\end{boundarybox}
